%------------------------------------------------------------------------------%
% Luis A. D'Afonseca
%
%------------------------------------------------------------------------------%
\documentclass[a4paper,12pt,fleqn]{article}

\usepackage{style_avaliacao}

\ShowAnswers

%------------------------------------------------------------------------------%
\begin{document}

\makeHeader{Integrais e Séries}{Prova 3 -- Turma 8}{14/01/2024}

\addGrowthOrder

% 01
%------------------------------------------------------------------------------%
\exercise{25}
Calcule o limite da sequência de somas parciais da série \;
\(
  \sum_{n=1}^\infty \frac{\pi}{n(n+1)}
\)

\begin{answer}
  Usando frações parciais, podemos reescrever o termo geral
  \[
    a_n
    = \frac{\pi}{n(n+1)}
    = \frac{\pi}{n} - \frac{\pi}{n+1}
  \]

  Expandindo a soma parcial \( S_N \), temos
  \[
    S_N
    = \sum_{n=1}^N \left( \frac{\pi}{n} - \frac{\pi}{n+1} \right)
    = \left( \frac{\pi}{1} - \frac{\pi}{2} \right)
    + \left( \frac{\pi}{2} - \frac{\pi}{3} \right)
    + \cdots
    + \left( \frac{\pi}{N} - \frac{\pi}{N+1} \right)
  \]
  Após o cancelamento telescópico, restam
  \[
    S_N = \pi - \frac{\pi}{N+1}
  \]
  Calculando o limite
  \[
    \lim_{N \to \infty} S_N
    = \lim_{N \to \infty} \pi - \frac{\pi}{N+1}
    = \pi
  \]
\end{answer}

% 02
%------------------------------------------------------------------------------%
\exercise{25}
Analise a convergência da série \;
\(
  \sum_{n=1}^\infty \frac{(n!)^2}{(2n)!}
\)

\begin{answer}
  \[
    a_n
    = \frac{(n!)^2}{(2n)!}
    = \frac{n!n!}{(2n)!}
  \]
  \[
    a_{n+1}
    = \frac{[(n+1)!]^2}{[2(n+1)]!}
    = \frac{(n+1)!(n+1)!}{(2n+2)!}
  \]
  \begin{align*}
    \frac{a_{n+1}}{a_n}
    & = \frac{(n+1)!(n+1)!}{(2n+2)!}\,\frac{(2n)!}{n!n!} \\[2mm]
    & = \frac{(n+1)\,n!\,(n+1)\,n!\,(2n)!}{(2n+2)(2n+1)(2n)!\,n!\,n!} \\[2mm]
    & = \frac{(n+1)(n+1)}{(2n+2)(2n+1)} \\[2mm]
    & = \frac{(n+1)(n+1)}{2(n+1)(2n+1)} \\[2mm]
    & = \frac{n+1}{4n+2} \\[2mm]
    & = \frac{1+\sfrac{1}{n}}{4+\sfrac{2}{n}}
  \end{align*}
  \[
    \rho
    = \lim_{n\to\infty} \frac{a_{n+1}}{a_n}
    = \lim_{n\to\infty} \frac{1+\sfrac{1}{n}}{4+\sfrac{2}{n}}
    = \frac{1 + 0}{4 + 0}
    = \frac{1}{4}
  \]
  Como \(\rho = \frac{1}{4} < 1\) a série converge pelo teste da razão.
\end{answer}

% 03
%------------------------------------------------------------------------------%
\exercise{25}
Verifique se a série converge \;
\(
  \sum_{n=1}^\infty \left[\frac{\sin(n)}{n}\right]^2
\)

\begin{answer}
  Como os termos da série
  \[
    \sum_{n=1}^\infty \left[\frac{\sin(n)}{n}\right]^2
    = \sum_{n=1}^\infty \frac{\sin^2(n)}{n^2}
  \]
  são não-negativos, podemos usar o teste da comparação.
  Observando que
  \(
    0 \leq \sin^2(n) \leq 1
  \)
  temos que
  \[
    a_n = \frac{\sin^2(n)}{n^2} \leq \frac{1}{n^2}
  \]
  Como \(b_n = \frac{1}{n^2}\) são os termos de uma $p$-série convergente
  (\(p=2>1\))
  o teste da comparação garante que a série
  \[
    \sum_{n=1}^\infty \frac{\sin^2(n)}{n^2}
  \]
  converge
\end{answer}

% 04
%------------------------------------------------------------------------------%
\exercise{25}
Use o teste da integral para mostrar que a soma da série existe \;
\(
  \sum_{n=1}^{\infty} \frac{1}{n(n+1)}
\)

\begin{answer}
  Vamos usar a função real
  \[
    f(x) = \frac{1}{x(x+1)}
  \]

  Precisamos verificar as três condições do teste
  \begin{enumerate}[label=\roman*)]
    \item \(f(x)>0\) sempre que \(x(x+1) > 0\).
          isso é, $x<-1$ ou $x>0$

    \item \(f(x)\) é contínua em para todo $x\in\R$, desde que $x\neq 0$
          e $x\neq -1$

    \item A derivada de \(f(x)\) é
    \begin{align*}
      f'(x)
      & = \left[\frac{1}{x(x+1)}\right]' \\[2mm]
      & = \left[\frac{1}{x^2+x}\right]' \\[2mm]
      & = \frac{0-(2x+1)}{x^2(x+1)^2} \\[2mm]
      & = -\frac{2x+1}{x^2(x+1)^2}
    \end{align*}
    A derivada será negativa sempre que \(2x+1>0\), ou seja,
    \(f(x)\) é decrescente quando \(x>-\frac{1}{2}\)
  \end{enumerate}

  Portanto, a função \(f(x) = \frac{1}{x(x+1)}\) é apropriada para
  o teste da integral. Vamos então calcular sua primitiva.
  \[
    F(x) = \int \frac{1}{x(x+1)} dx
  \]
  Por frações parciais temos que
  \[
    \frac{1}{x(x+1)} = \frac{1}{x} - \frac{1}{x+1}
  \]
  então
  \[
    F(x)
    = \int \frac{1}{x} - \frac{1}{x+1} dx
    = \ln\abs{x} - \ln\abs{x+1} + C
  \]
  Calculando a integral imprópria
  \begin{align*}
    A
    & = \int_1^\infty f(x) dx \\[2mm]
    & = \lim_{b\to\infty}\int_1^b f(x)dx \\[2mm]
    & = \lim_{b\to\infty} F(x)\evalat{1}{b} \\[2mm]
    & = \lim_{b\to\infty} \big(\ln\abs{x} - \ln\abs{x+1}\big)\evalat{1}{b} \\[2mm]
    & = \lim_{b\to\infty}
        \Big(
          \big( \ln\abs{b} - \ln\abs{b+1}\big) -
          \big( \ln\abs{1} - \ln\abs{1+1}\big)
        \Big) \\[2mm]
    & = \lim_{b\to\infty} \left( \ln\abs{\frac{b}{b+1}}\right)
      + \ln(2) \\[2mm]
    & = \ln\abs{\lim_{b\to\infty} \frac{b}{b+1}}
      + \ln(2) \\[2mm]
    & = \ln(1)
      + \ln(2) \\[2mm]
    & = \ln(2)
  \end{align*}
  Como a integral converge a série também converge.
\end{answer}

%% 05
%%------------------------------------------------------------------------------%
%\exercise{20}
%Determine quantos termos precisam ser somados na série do exercício anterior
%para que obter uma precisão de pelo menos \(10^{-6}\).
%
%\begin{answer}
%  Sabemos que
%  \[
%    \abs{R_n} \leq \int_n^\infty \frac{1}{x(x+1)} dx
%  \]
%
%  Calculando a integral imprópria
%  \begin{align*}
%    A
%    & = \int_n^\infty f(x) dx \\[2mm]
%    & = \lim_{b\to\infty} \big(\ln\abs{x} - \ln\abs{x+1}\big)\evalat{n}{b} \\[2mm]
%    & = \lim_{b\to\infty}
%      \Big(
%      \big( \ln\abs{b} - \ln\abs{b+1}\big) -
%      \big( \ln\abs{n} - \ln\abs{n+1}\big)
%      \Big) \\[2mm]
%    & = \ln\abs{\lim_{b\to\infty} \frac{b}{b+1}}
%      - \ln\left(\frac{n}{n+1}\right) \\[2mm]
%    & = - \ln\left(\frac{n}{n+1}\right) \\[2mm]
%    & = \ln\left(\frac{n+1}{n}\right)
%  \end{align*}
%  Como a integral converge a série também converge.
%\end{answer}
%------------------------------------------------------------------------------%
\end{document}
%------------------------------------------------------------------------------%
